% TOP_PROBLEM: {"id": 2560, "title": "Kourovka Notebook Problem 21.51", "category_id": 17, "category": {"id": 17, "name": "group_theory", "display_name": "Group Theory", "description": "Problems about groups, group actions, representations, and related algebraic structures.", "slug": "group-theory", "order_index": 17, "created_at": "2026-07-31T15:26:25.671Z"}, "status": "open", "problem_number": "KOU-21.51", "set_id": 10}
% FIRST_POSTED: 2026-10-01
% ENTRY_KIND: statement_only
% UnsolvedMath ID 2560; source catalogue code KOU-21.51
% !TeX program = lualatex
% Definition and sources only; this file is not an LLM solution attempt.
\documentclass[11pt,a4paper]{article}
\usepackage[margin=25mm]{geometry}
\usepackage{fontspec}
\setmainfont{lmroman10-regular.otf}[BoldFont=lmroman10-bold.otf,
 ItalicFont=lmroman10-italic.otf,BoldItalicFont=lmroman10-bolditalic.otf]
\usepackage{amsmath,amssymb,mathtools}
\usepackage{unicode-math}
\setmathfont{latinmodern-math.otf}
\AtBeginDocument{\let\setminus\smallsetminus}
\usepackage{xurl,hyperref}
\hypersetup{colorlinks=true,urlcolor=blue}
\setcounter{secnumdepth}{0}
\setlength{\parindent}{0pt}
\setlength{\parskip}{5pt}
\setlength{\emergencystretch}{3em}
\begin{document}
\section*{Kourovka Notebook Problem 21.51}\label{2560}
{\small UnsolvedMath ID 2560 · KOU-21.51 · Group Theory · Kourovka Notebook - New Problems, Issue 21}
\subsection{Definitions and mathematical statement}
Fix a prime $p$ and a positive integer $n$. A finite $p$-group is a
group of order $p^a$ for some integer $a\ge0$. An abelian subgroup
has commuting elements. An elementary abelian $p$-subgroup is a subgroup
isomorphic to $(\mathbb Z/p\mathbb Z)^r$ for some $r$, or equivalently
an abelian subgroup in which every nonidentity element has order $p$.
Normality of a subgroup means invariance under conjugation by every
element of the ambient group.

Part (a) asks for precisely those pairs $(p,n)$ with the property
\[
 \forall P\text{ finite }p\text{-group},\quad
 [\exists A\le P\text{ abelian},\ |A|=p^n]
 \Longrightarrow
 [\exists B\trianglelefteq P\text{ abelian},\ |B|=p^n].
\]
Part (b) asks for the analogous classification when both occurrences
of ``abelian'' are replaced by ``elementary abelian.''
For a fixed prime, this is a determination of all positive integers
$n$ for which the corresponding implication is universal in $P$.

The normal subgroup supplied in the conclusion need not contain the
original subgroup, be contained in it, or be conjugate to it. It must
have exactly the same order, not merely some smaller nontrivial order.
No bound is imposed on $|P|$ beyond finiteness and its prime-power
form. A negative answer for a pair requires a $p$-group with the
specified subgroup but no normal subgroup of that type and order.
The two classifications are separate requirements.
\subsection{Short English statement}
For which primes and subgroup orders does existence of an abelian, or elementary abelian, subgroup force a normal subgroup of the same type and order?
\subsection{Sources}
\begin{itemize}
\item[S1] ULAM.AI and the UnsolvedMath Contributors, supplied problems.json, record 2560 (KOU-21.51), Kourovka Notebook - New Problems, Issue 21. Catalogue identity and source transcription; CC BY 4.0.
\url{https://huggingface.co/datasets/ulamai/UnsolvedMath}.
\item[S2] The Kourovka Notebook, 21st edition, new problems, Problem 21.51. Expanded formulation of the supplied catalogue statement.
\url{https://arxiv.org/pdf/1401.0300v47}.
\end{itemize}
\end{document}
